% Standalone source for the routed-gate circuit figure.
%
% Build:   pdflatex circuit_routed.tex      (from paper_figures/)
% Include: \includegraphics{circuit_routed}     <-- NO width= option
%
% The output PDF is built at exactly 246.0pt wide, which is \columnwidth for
% main.tex's class options (aip,reprint,jcp / revtex4-1).  Included at its
% natural size it therefore fills the column exactly, and nothing can ever
% rescale it.  Verified: TeX reads the natural width as 245.99942pt.
%
% Do NOT use \includegraphics[width=\columnwidth]{...} here.  It is scale 1.0
% in the normal in-column case, but if the float ever gets deferred to a float
% page, revtex typesets that page in one-column mode where \columnwidth is
% \textwidth = 510pt -- and the figure silently comes out 2.07x too big, with
% labels dwarfing the body text.  Omitting width= is immune to this.
%
% If you change the class options, re-check \the\columnwidth and update the
% 246.0pt below, otherwise the figure stops matching the column.
%
% (pdfinfo reports 245.081pt because PDF uses big points, 72bp/in, whereas TeX
% pt is 72.27/in: 246 * 72/72.27 = 245.081.  The page really is 246.0pt.)
%
% Font setup below mirrors main.tex (T1 + mathptmx, 10pt base, so \footnotesize
% is 8pt in both).  Keep them in sync or the labels will not match the paper.
%
% The (a)/(b) panel labels are baked in here on purpose: that way main.tex does
% NOT need the subcaption package.  subcaption pulls in caption, which does not
% recognise revtex4-1 and silently re-formats every caption in the paper from
% justified to centred.
%
% Layout constants (row sep, column sep, the \smash/\raisebox around the gate
% label) are explained in circuits.tex -- change them there and copy across.

\documentclass{standalone}
\usepackage[T1]{fontenc}
\usepackage{mathptmx}
\usepackage{amsmath}
\usepackage{graphicx}
\usepackage{tikz}
\usepackage{quantikz}

\newcommand{\isw}{\raisebox{-0.5\height}[1.5pt][1.5pt]{\rotatebox{90}{\mbox{\footnotesize$\sqrt{i\text{SWAP}}$}}}}

% Panel labels pinned to 9pt/10.5pt: that is revtex4-1's native caption font
% (measured), so these match the \caption text of the figure they sit in.  Set
% explicitly rather than via \small etc. so they cannot drift if the class or
% its base size changes -- this file is compiled by `standalone`, whose
% relative size commands need not agree with revtex's.
%
% Note this is deliberately NOT the same as the 8pt \footnotesize used for the
% \sqrt{iSWAP} gate labels inside \isw above; do not "unify" them.  The gate
% label size feeds the geometry documented in circuits.tex (its 30.15pt
% rotated height sets the minimum row sep), so changing it breaks the layout.
\newcommand{\panellabel}[1]{{\fontsize{9}{10.5}\selectfont #1}}

\newsavebox{\seqbox}\newsavebox{\parcircbox}

\begin{document}
% Each circuit is typeset into a box first and its panel is made exactly that
% wide.  Fixed-fraction panel widths let a circuit overflow its slot, and
% standalone then crops whatever sticks out past 246pt (this silently cut the
% last SWAP off the sequential circuit).  \hfill splits the leftover space.
\begin{lrbox}{\seqbox}
\begin{quantikz}[row sep={0.85cm,between origins}, column sep=0.3cm]
\lstick{$q_0$} & \qw      & \qw      & \gate[wires=2]{\isw} & \qw      & \qw      & \qw \\
\lstick{$q_1$} & \qw      & \swap{1} &                      & \swap{1} & \qw      & \qw \\
\lstick{$q_2$} & \swap{1} & \targX{} & \qw                  & \targX{} & \swap{1} & \qw \\
\lstick{$q_3$} & \targX{} & \qw      & \qw                  & \qw      & \targX{} & \qw
\end{quantikz}
\end{lrbox}%
\begin{lrbox}{\parcircbox}
\begin{quantikz}[row sep={0.85cm,between origins}, column sep=0.3cm]
\lstick{$q_0$} & \swap{1} & \qw                  & \swap{1} & \qw \\
\lstick{$q_1$} & \targX{} & \gate[wires=2]{\isw} & \targX{} & \qw \\
\lstick{$q_2$} & \swap{1} &                      & \swap{1} & \qw \\
\lstick{$q_3$} & \targX{} & \qw                  & \targX{} & \qw
\end{quantikz}
\end{lrbox}%
\begin{minipage}{246.0pt}
\begin{minipage}[b]{\wd\seqbox}
\centering
\usebox{\seqbox}

\vspace{2ex}
\panellabel{(a) Serial}
\end{minipage}
\hfill
\begin{minipage}[b]{\wd\parcircbox}
\centering
\usebox{\parcircbox}

\vspace{2ex}
\panellabel{(b) Parallel}
\end{minipage}
\end{minipage}
\end{document}
